(1) 相速度 $v_p=\omega/k=c/n(\omega)$。群速度为 $v_g=[dk/d\omega]^{-1}=c/[n(\omega)+\omega\,dn/d\omega]$，其分母定义为群折射率 $n_g=n+\omega n'$；载波相位前进与脉冲包络传输一般不能用同一速度描述。
(2) 在中心频率处传播相位为 $\phi_0=k(\omega_0)L=n(\omega_0)\omega_0L/c$，群时延 $\tau_g=d[k(\omega)L]/d\omega|_{\omega_0}=Ln_g(\omega_0)/c$。此处相位延迟时间 $Ln(\omega_0)/c$ 与群时延相差 $L\omega_0n'(\omega_0)/c$。
(3) 定义二阶谱相位系数 $\beta_2=k''(\omega_0)$，传播时频谱增加 $\exp[i\beta_2 L\Omega^2/2]$。高斯积分的复宽参数为 $\tau_0^2-i\beta_2L/2$（符号取决于Fourier约定），强度脉宽标准差为 $\sigma_t(L)=\tau_0\sqrt{1+[\beta_2L/(2\tau_0^2)]^2}$，且脉冲中心在 $\tau_g$ 处。
(4) 若 $n$ 与频率无关，则 $n_g=n$、$\beta_2=0$，高斯脉冲仅整体延迟而不展宽。在强吸收或反常色散情况下群速度可能不能直接代表信息前沿速度；二阶模型要求带宽 $\Delta\omega\sim1/\tau_0$ 足够窄，且传播长度不使三阶及以上相位项显著。
